% Figure: pigeonhole for the 51-numbers example --- odd-core classes.
% Source for Figure~\ref{fig:pigeon}. TikZ only; regenerable by compiling main.tex.
\begin{tikzpicture}[
  box/.style={draw=rulegray,rounded corners=2pt,minimum height=8mm,inner sep=3pt},
  lab/.style={font=\footnotesize\sffamily,text=slate},
]
% 5 of the 50 classes: odd part 1,3,5,7,9 ... 99
\node[box,align=center,text width=1.55cm] at (0,0) {{$1,2,4,8,$\\$16,32,64$}};
\node[box,align=center,text width=1.55cm] at (2.1,0) {{$3,6,12,$\\$24,48,96$}};
\node[box,align=center,text width=1.55cm] at (4.2,0) {{$5,10,20,$\\$40,80$}};
\node[box,align=center,text width=1.55cm] at (6.3,0) {{$7,14,28,$\\$56$}};
\node[lab] at (7.9,0) {\dots};
\node[box,align=center,text width=1.55cm] at (9.4,0) {$99$};
\node[lab] at (0,0.9) {odd part $=1$};
\node[lab] at (2.1,0.9) {odd part $=3$};
\node[lab] at (4.2,0.9) {odd part $=5$};
\node[lab] at (6.3,0.9) {odd part $=7$};
\node[lab] at (9.4,0.9) {odd part $=99$};
\node[lab] at (4.7,-0.9) {50 classes cover $1..100$; two picks in one class differ by a factor of 2};
\node[lab] at (4.7,-1.4) {the all-odd 50-set shows 50 is the true threshold};
\end{tikzpicture}
